\documentclass[10pt]{article}

\def\Type{LessonDJ}
\def\TypeNum{Workshop Week 14} %Lesson #
\def\TypeTitle{Review for Final Exam} %Lesson Title
\def\Semester{Fall 2021} %Not Used 
\def\Course{MATH 1190}%Course number and name
\def\Targets{   FE1 through FE 10 }

\usepackage{preambleDJ} 

%%%%%%%%%%%% BEGIN DOCUMENT %%%%%%%%%%%%%%%
%%%%%%%%%%%% BEGIN DOCUMENT %%%%%%%%%%%%%%%
%%%%%%%%%%%% BEGIN DOCUMENT %%%%%%%%%%%%%%%
%%%%%%%%%%%% BEGIN DOCUMENT %%%%%%%%%%%%%%%
%%%%%%%%%%%% BEGIN DOCUMENT %%%%%%%%%%%%%%%

\begin{document}
\pagestyle{otherpages}
\thispagestyle{firstpage}

\vs{.6}
\section{Derivatives}
Find the derivative of the following functions:
\begin{enumerate}
    \item $ 7x^4-3x^2+10x-2$
    \item $\ds \frac{1}{x^3}+2\sqrt{x}$
    \item $(3x^2-5)(2x+1)$
    \item $\ds \frac{x^2+1}{x-4}$
    \item $\sqrt{3x^2+9x-7}$
    \item $(8x^2-1)^7$
    \item $ e^{3x^2-x}$
    \item $x\ln(4x^2+1)$
    
\end{enumerate}
\newpage
\section{Definite Integrals}
Find the following integrals:
\begin{enumerate}
    \item $\ds \int_0^1\ 3x^2+4x-1\ dx$
    \item $\ds \int  (2x^2+1)(6x^3+2x)\ dx$
    \item $\ds \int_{-1}^1\ 2e^x+\log_7(2)\ dx$
    \item $\ds \int_1^2\ \frac{1}{x^3}+\frac{2}{x^2}+\frac{3}{x}\ dx$
    \item $\ds \int_1^2\ \frac{x^2-1}{x-1}\ dx$
    \item $\ds \int \ \frac{3\sqrt{z^3}+2z+1}{z^2}\ dz$
    \item $\ds \int (x^2+1)(\frac{1}{x}+\frac{1}{x^2})\ dx$
    \item $\ds\int_0^2\ 4\cdot5^x\ dx$
    
\end{enumerate}

\newpage
\section{Optimization}
\begin{enumerate}
\item A rectangle is inscribed under the parabola $y=16-x^2$ with its base on the x-axis. Now we want to find the maximum possible area of the rectangle.

{\bf Modeling Phase} We first need to develop a formula for the area of such a rectangle and reduce it to an one variable function.
{}
\tasksA{2}{
\task Draw a graph of $y=16-x^2$ with a rectangle inside of it.  The base of the rectangle will be on the $x$-axis and the rectangle will extend up until it touches the graph of $y=16-x^2$.
\task How can the usual variables $x$ and $y$ be used to describe the height and width of the rectangle? Use these variables to create a function that describes the area of the rectangle. \vs{1.5}


\task What is the possible range of widths of the rectangle? How does this relate to the possible values for $x$?
\task Use the formula for the parabola $y=16-x^2$ to reduce the area formula into an one variable function.}
\vs{1}
{\bf Calculus Phase} Think about the process you have been using to optimize quantities. Use that process to find the maximal area. Don't forget to check that your final answer of the width is indeed within the range of the possible width in Part (c).


    \vfill
    \end{enumerate}

\newpage
\section{Summations}
\begin{enumerate}

    
    \item Rewrite in summation notation -you don't need to solve!-:
    \begin{enumerate}
        \item $\ds 3 + 9 + 27 + 81 $
        \item $\ds 3 - 9 + 27 - 81 + 243$
        \item $\ds -\frac{5}{2}+\frac{10}{4}-\frac{15}{8}+\frac{20}{16}$


    \end{enumerate}

    \item Knowing that $\ds \sum_{i=1}^{n}a_i=6,\sum_{i=1}^{n}b_i=7,\sum_{i=1}^{n}c_i=-5$, use the properties of summation to compute the values of the following summations:
    \begin{enumerate}
        \item $\ds \sum_{i=1}^n \frac{b_i-a_i}{9}$
        \item $\ds \sum_{i=1}^n a_i-2b_i+3c_i-4$\hspace{10mm} (Your answer should depend on $n$.)
    \end{enumerate}
        \item Compute explicitly the following sums: 
    \begin{enumerate}
        \item $\ds \sum_{k = 1}^4 (-1)^k\cdot 2k$
        \item $\ds \sum_{k = 0}^3 (-1)^{k+1}\cdot 2(k+1)$
        \item $\ds \sum_{k = 1000}^{1004} (-1)^{k-999} \cdot2(k-999)$
    \end{enumerate}
    \end{enumerate}




\newpage
\section{Probability (Discrete)}
\begin{enumerate}
%    \item 
%In a tabletop game, a player rolls two fair 4-sided dice $(D_1,D_2)$. Let the random variable $X$ represent the absolute difference $|D_1 - D_2|$ between the numbers rolled on the two dice.
%\begin{enumerate}
%    \item[(a)] List out the sample space $\Omega$ for rolling the two 4-sided dice. 
%    \item[(b)] Determine the probability mass function for $X$ and present it in a table.
%    \item[(c)] Draw a probability histogram for the random variable $X$.
%    \item[(d)] Compute the expected value $E[X]$.
%    \item[(e)] Compute the variance $Var[X]$.
%    \item[(f)] What is the probability that the absolute difference is an even number? (Note that $0$ is an even number)
%\end{enumerate}
%\vspace{4cm}
\item At a local fundraiser, attendees can only participate in a different raffle game. The player draws a two tokens from a box containing contains $3$ tokens. One is worth $\$2$, another is worth $\$4$, and the last token is worth $\$6$. After the first draw, the token is put back into the box.\\\\
Let the random variable $X$ represent the total sum of the prize money won after two draws.
\enum{
\item Write out all possible combinations that can be drawn and their sums. {\em Hint}: It might be helpful to create a table similar to the one when we examined the sum of $2$ dice.
\item What is the probability that a participant will wind $\$10$ or more?
\item On average, how much will a participant win?
\item What is the variance, $Var(X)$?
}
% \item In a charity raffle, a contestant draws exactly one token from each of three boxes:
% \begin{itemize}
%     \item Box A contains $2$ tokens worth $\$1$ and $\$5$.
%     \item Box B contains $3$ tokens worth $\$1$, $\$2$, and $\$3$.
%     \item Box C contains $2$ tokens worth $\$3$ and $\$5$
% \end{itemize}
% Assume all tokens within a given box are equally likely to be drawn. Let the random variable $X$ represent the total sum of the prize money drawn from all three boxes.
% \begin{enumerate}
% \item[(a)] List out the sample space $\Omega$ for the tokens drawn from the three boxes.
% \item[(b)] Determine the probability mass function for $X$.
% \item[(c)] Draw a probability histogram for the random variable $X$.
% \item[(d)] Compute the expected value $E[X]$.
% \item[(e)] Compute the variance $Var[X]$.
% \item[(f)] What is the probability that the total prize is at least $\$10$.
% \end{enumerate}
\end{enumerate}

\newpage
\section{Probability (Continuous)}
Let $X$ be a continuous random variable representing the lifespan (in thousands of hours) of a light bulb. The probability density function of $X$ is given by
$$f(x) = \begin{cases} 
c x (2 - x) & \text{for } 0 \le x \le 2 \\ 
0 & \text{otherwise} 
\end{cases}$$
\begin{enumerate}[]
    \item[(a)] Find the value of the constant $c$ that makes $f(x)$ a valid probability density function.
    \item[(b)] Find the cumulative distribution function $F(x)$ for this random variable.
    \item[(c)] Compute the expected lifespan $E[X]$ of that light bulb.
    \item[(d)] Compute the variance $Var[X]$.
    \item[(e)] What is the probability that a light bulb lasts longer than 1,500 hours?
\end{enumerate}

\end{document}